\section{Begrænsninger og eksterne signaler}

\subsection{P-Lead-Lag og actuatorbegrænsning / Lag, saturation, windup}
\label{sec:limits}

\paragraph{Brug når}
Der angives Lag-parameter \(\beta\), rate/amplitudelimit eller integrator windup.

\paragraph{Input fra opgaven}
For Lag: \(\omega_c,N_i,\alpha,\gamma_M\) og plantfase. For limits:
\(u_{\min},u_{\max}\) eller \(R_{\max}\), reference og controllerrespons.

\paragraph{Metode}
\begin{enumerate}
\item Afgør typen ved at se på elementet: et Lag-led er en transferfunktion med
      \(\beta\) og behandles lineært med \eqref{eq:lag-controller-phase}; saturation
      eller rate limit er en ulighed på \(u(t)\) eller \(\dot u(t)\) og behandles med
      \eqref{eq:actuator-limits}.
\item Ved Lag-design: skriv først den fase som allerede kommer fra plant, P/PI og Lead
      ved \(\omega_c\). Den manglende fase er
      \(\phi_{\mathrm{lag}}=-180^\circ+\gamma_M-\phi_{\mathrm{uden\ lag}}\), og den
      skal være negativ. Indsæt \(\phi_{\mathrm{lag}}\) og \(N_i\) i
      \eqref{eq:lag-controller-phase}, og løs for \(\beta\); accepter kun \(\beta>1\).
\item Ved amplitudesaturation: beregn eller skitser først den ubegrænsede
      controllerkommando \(u_{\mathrm{lin}}(t)\). Hvis maksimum er større end
      \(u_{\max}\), eller minimum er mindre end \(u_{\min}\), erstattes signalet af
      grænsen i \eqref{eq:actuator-limits}, og den lineære model beskriver ikke længere
      hele responsen.
\item Ved rate limit: beregn den stejleste ændring i kommandoen, enten som
      \(\max|\dot u_{\mathrm{lin}}(t)|\) eller fra en tangent på plottet. Hvis den er
      større end \(R_{\max}\), kan aktuatoren ikke følge den lineære kommando.
\item Windup mistænkes når en integrator fortsætter med at akkumulere fejl, mens
      aktuatoren står på en grænse. Vælg mitigation ud fra årsagen: lavere bandwidth
      reducerer \(u\)-kravet, prefilter reducerer referencespring, integrator limit eller
      anti-windup stopper integratoropbygning, og Lag ændrer lavfrekvent/højfrekvent
      gain uden at være en hard limit.
\end{enumerate}

\paragraph{Formler}
\begin{equation}
\label{eq:lag-controller-phase}
C_{\mathrm{lag}}(s)=\frac{\tau_i s+1}{(\tau_i/\beta)s+1},\quad \beta>1,\qquad
\phi_{\mathrm{lag}}=\arctan\!\left(\frac{N_i(1-\beta)}{1+\beta N_i^2}\right),
\end{equation}
\begin{equation}
\label{eq:actuator-limits}
u_{\min}\leq u_{\mathrm{sat}}(t)\leq u_{\max},\qquad |\dot u(t)|\leq R_{\max}.
\end{equation}

\paragraph{Symboler}
\(\beta\) er lag-forhold; \(N_i=\omega_c\tau_i\); \(R_{\max}\) er maksimal inputrate.

\paragraph{Antagelser}
Lag er lineært og giver ikke en ren integrator. Saturation og rate limit er
ikke-lineære, så en rent lineær marginanalyse er ikke tilstrækkelig for store signaler.

\paragraph{Hurtigtjek}
\(\phi_{\mathrm{lag}}<0\) og \(\beta>1\). Ved windup ses typisk lang recovery efter saturation.
En rate limiter er ikke det samme som amplitudesaturation: amplituden kan være inden for
grænsen, selv om \(|\dot u|\) er for stor. I slides' limited-system-design vælges
crossover ofte tæt på eller lidt over den dominerende pol, men designet skal simuleres
eller vurderes med rate limiteren slået til.

\paragraph{Typiske fælder}
At behandle Lag som PI med garanteret nul stepfejl; at erklære stor-step-stabilitet
ud fra en small-signal transferfunktion.

\paragraph{Python}
\code{solve\_lag\_beta(lag\_phase\_deg,n\_i)} beregner \(\beta\) fra den fase,
du selv har udledt (F21 Q17). Der automatiseres ikke limit- eller windupdesign.

\subsection{Disturbance og sensitivity / \(S\), \(T\), prefilter}
\label{sec:disturbance}

\paragraph{Brug når}
En disturbance, measurement noise eller et prefilter indgår, og effekt på output/fejl
skal bestemmes.

\paragraph{Input fra opgaven}
Disturbancens indsættelsespunkt og fortegn, loop \(L\), ønsket outputsignal og eventuelt
stepstørrelse.

\paragraph{Metode}
\begin{enumerate}
\item Vælg ét uafhængigt input ad gangen. Hvis du finder \(Y/D\), sæt \(R=0\) og
      \(N=0\); hvis du finder \(Y/R\), sæt disturbances og noise til nul.
\item Marker hvor disturbance kommer ind. Kommer den direkte på outputtet efter
      planten, bruger du output-disturbance-rækken i \eqref{eq:disturbance-standard-paths};
      kommer den ind før planten, bruger du plant-input-rækken; kommer den ind som
      målefejl i feedbacksignalet, bruger du noise-rækken.
\item Hvis diagrammet er standard negativ unity feedback, skriv først \(L=CG\) og brug
      \(S,T\) fra \eqref{eq:sensitivity-functions}. Hvis der er sensor \(H\), ekstra
      feed-forward eller anderledes fortegn, skriv blokdiagramligninger som i
      \S\ref{sec:feedback} og brug kun \(S/T\)-navnene efter at formen matcher.
\item Ved constant disturbance: sæt \(D(s)=d_0/s\), find \(Y(s)\) eller \(E(s)\), og
      brug final value theorem \eqref{eq:steady-state-error} efter stabilitetscheck.
\item For et referenceprefilter: gang kun referencevejen med \(F(s)\) som i
      \eqref{eq:reference-prefilter}. Fordi prefilteret står uden for feedbacksløjfen,
      ændrer det ikke \(L\), phase margin, gain margin eller disturbance rejection.
\end{enumerate}

\paragraph{Formler}
\begin{equation}
\label{eq:sensitivity-functions}
S(s)=\frac1{1+L(s)},\qquad T(s)=\frac{L(s)}{1+L(s)},\qquad S+T=1,
\end{equation}
\begin{equation}
\label{eq:disturbance-standard-paths}
\text{output disturbance: } \frac{Y}{D}=S,\quad
\text{plant-input disturbance: } \frac{Y}{D}=GS,\quad
\text{measurement noise: } \frac{Y}{N}=-T,
\end{equation}
\begin{equation}
\label{eq:reference-prefilter}
\text{referenceprefilter: }\frac{Y}{R}=F(s)T(s).
\end{equation}

\paragraph{Symboler}
\(D,N\) er disturbance/noise; \(F\) er referenceprefilter; \(S,T\) er
dimensionsløse standard-sensitiviteter.

\paragraph{Antagelser}
De kompakte \(S/T\)-former antager standard negativ unity feedback og den angivne
placering; et konkret blokdiagram har forrang.

\paragraph{Hurtigtjek}
Et prefilter uden for loopet må ikke ændre \(L\), PM/GM eller disturbance rejection.
Stor lavfrekvent loopgain reducerer typisk outputdisturbance via \(S(0)\).
Ved flere disturbancekurver bruges DC-niveauet først: en disturbance efter en integrator
kan have nul steady-state-output, mens en disturbance før integratoren kan kræve et
ikke-nul controllerinput og give et lavfrekvent niveau bestemt af \(K_P\).

\paragraph{Typiske fælder}
At bruge fejlsignalet fra en korrupt måling som performancefejl eller bytte input- og
outputdisturbance.

\paragraph{Python}
Ingen offentlig funktion reducerer disturbanceblokdiagrammer; signalvej og fortegn er
eksamensmetoden.

\subsection{Dynamisk feed-forward / measured disturbance cancellation}
\label{sec:feedforward}

\paragraph{Brug når}
En målbar disturbance \(d\) føres gennem \(D(s)\), og en feed-forward-blok \(F_d(s)\)
skal afvise den gennem plant path \(G(s)\).

\paragraph{Input fra opgaven}
\(G(s)\), \(D(s)\), summationsfortegnet \(\sigma_D\in\{-1,+1\}\), og om disturbance
kan måles.

\paragraph{Metode}
\begin{enumerate}
\item Skriv fortegnet \(\sigma_D\) ud fra summationspunktet: \(+1\) hvis disturbance
      lægges til outputbidraget, \(-1\) hvis den trækkes fra. Skriv derefter
      disturbancebidraget før loopnævneren som i tælleren i
      \eqref{eq:feedforward-cancellation}.
\item Sæt disturbance-numeratoren lig nul:
      \(\sigma_DD(s)+G(s)F_d(s)=0\). Isoler \(F_d\), hvilket giver første del af
      \eqref{eq:feedforward-cancellation}.
\item Kontrollér properness ved at sammenligne polynomiumsgrader i \(F_d\): nævnergrad
      skal være mindst tællergrad for en proper realisering. Kontrollér stabilitet ved
      at finde polerne i \(F_d\); de skal ligge strengt i \(\LHP\).
\item Hvis den ideelle \(F_d\) er improper eller ustabil, må perfekt dynamisk
      cancellation ikke loves. Brug i stedet en statisk approximation ved lav frekvens,
      \(F_d(0)\), hvis den findes, eller multiplicér med et stabilt lavpasfilter og
      skriv at rejection bliver begrænset.
\item Indsæt den valgte \(F_d\) tilbage i \eqref{eq:feedforward-cancellation}. Hvis
      numeratoren bliver nul, er rejection perfekt i nominalmodellen; hvis ikke, er
      restleddet den disturbance der stadig når outputtet.
\end{enumerate}

\paragraph{Formler}
\begin{equation}
\label{eq:feedforward-cancellation}
F_d(s)=-\frac{\sigma_DD(s)}{G(s)},\qquad
G_{yd}(s)=\frac{\sigma_DD(s)+G(s)F_d(s)}{1+C(s)G(s)}.
\end{equation}

\paragraph{Symboler}
\(F_d\) er disturbance feed-forward; \(\sigma_D\) er fortegnet ved disturbance-summeringen.

\paragraph{Antagelser}
Disturbance er målt, plant og disturbance path er nominalt korrekte, og inversen kan
realiseres.

\paragraph{Hurtigtjek}
Indsæt \(F_d\) tilbage i numeratoren og se at den bliver nul; tjek poler og relative grader.

\paragraph{Typiske fælder}
Forkert fortegn, invertering af forkert plant path eller påstand om perfekt rejection ved
modelusikkerhed.

\paragraph{Python}
\code{ideal\_disturbance\_feedforward(...,disturbance\_sign)} danner ratioen og
rapporterer proper/stable efter at du har valgt fortegnet. For F21 Q20 er
\(\sigma_D=-1\), og nominalt fås \(G_{yd}=0\).

\subsection{Reference feed-forward for nul stationær fejl / inverse with low-pass}
\label{sec:tracking-feedforward}

\paragraph{Brug når}
Opgaven viser en reference feed-forward-blok \(F(s)\) omkring et feedbackloop og spørger,
hvilken \(F(s)\) der giver nul steady-state tracking error.

\paragraph{Input fra opgaven}
Plant \(G(s)\), feedbackgain/controller \(K_P\), krav om steady-state error, og om
plantinversen er proper/stabil.

\paragraph{Metode}
\begin{enumerate}
\item Skriv reference-to-output pathen ved lav frekvens. For unit step er kravet
      \(y_{\mathrm{ss}}=r_{\mathrm{ss}}\), altså DC-gain fra \(R\) til \(Y\) lig 1.
\item Den ideelle feed-forward til tracking er plantinversen \(F(s)=1/G(s)\), fordi
      \(G(s)F(s)=1\). Dette fjerner ikke behovet for feedback; det former referencevejen.
\item Kontrollér realiserbarhed. Hvis \(1/G(s)\) er improper, multiplicér med et stabilt
      lavpasfilter som i \eqref{eq:tracking-feedforward}. Vælg \(\tau_f\) lille nok til
      at lavfrekvent gain stadig er omtrent 1.
\item Brug ikke \(G(0)\) som feed-forward, hvis formålet er at invertere planten; den
      rigtige statiske invers ville være \(1/G(0)\), og den dynamiske invers er
      \(1/G(s)\) plus eventuelt lavpasfilter.
\end{enumerate}

\paragraph{Formler}
\begin{equation}
\label{eq:tracking-feedforward}
F_{\mathrm{track}}(s)=\frac{1}{G(s)(\tau_f s+1)},\qquad
G(s)F_{\mathrm{track}}(s)=\frac{1}{\tau_f s+1}.
\end{equation}

\paragraph{Symboler}
\(\tau_f\) er lavpasfilterets tidskonstant; \(F_{\mathrm{track}}\) er reference
feed-forward til tracking.

\paragraph{Antagelser}
Plantmodellen er nominalt korrekt, og lavpasfilteret er stabilt. Perfekt tracking gælder
kun lavfrekvent/asymptotisk, ikke nødvendigvis i transienten.

\paragraph{Hurtigtjek}
Ved \(s=0\) giver \eqref{eq:tracking-feedforward} \(G(0)F_{\mathrm{track}}(0)=1\), så
unit-step steady-state error kan blive nul.
Når plantinversen har højere tællergrad end nævnergrad, skal den gøres proper med et
stabilt lavpasfilter af tilstrækkelig orden. Vælg filtertidskonstanten lille i forhold
til de relevante planttidskonstanter, hvis opgaven kræver lavfrekvent tracking eller
disturbance cancellation.
